\documentclass[runningheads]{llncs}

% --------------------------------------------------------------- 
\usepackage[mobile]{eccv}


% ---------------------------------------------------------------
% Other packages
\newcommand{\ours}{WorldAgents}
% Commonly used abbreviations (\eg, \ie, \etc, \cf, \etal, etc.)
\usepackage{eccvabbrv}

% Include other packages here, before hyperref.
\usepackage{graphicx}
\usepackage{booktabs}

% The "axessiblity" package can be found at: https://ctan.org/pkg/axessibility?lang=en
\usepackage[accsupp]{axessibility}  % Improves PDF readability for those with disabilities.

% Reduce spacing around floats (ECCV-safe)
\setlength{\textfloatsep}{4pt}     % Space between floats and text
\setlength{\abovecaptionskip}{4pt} % Space above caption
\setlength{\belowcaptionskip}{4pt} % Space below caption
\setlength{\floatsep}{4pt}
\setlength{\intextsep}{4pt}
% for equations
\setlength{\abovedisplayskip}{6pt}
\setlength{\belowdisplayskip}{6pt}
\setlength{\abovedisplayshortskip}{4pt}
\setlength{\belowdisplayshortskip}{4pt}
\usepackage{inconsolata} % Adds a highly readable monospace font
\usepackage{listings}
\usepackage{xcolor}

% Define a clean style for your markdown/text listing
\lstdefinestyle{mdStyle}{
    backgroundcolor=\color{gray!5},   % light gray background
    basicstyle=\ttfamily\scriptsize,  % Use a smaller monospace font to save space
    breakatwhitespace=false,         
    breaklines=true,                  % Automatically wrap long lines
    captionpos=b,                     % Position caption at the bottom
    keepspaces=true,                 
    numbers=none,                     % No line numbers (change to 'left' if you want them)
    showspaces=false,                
    showstringspaces=false,
    showtabs=false,                  
    frame=single,                     % Adds a border around the text
    rulecolor=\color{gray!40},         % Color of the border
    breakindent=0pt,                  % <--- FIX: Removes the 20pt indent for wrapped lines
    breakautoindent=false,            % <--- FIX: Stops it from copying the previous line's indent
}

% ---------------------------------------------------------------
% Hyperref package

\usepackage{hyperref}

% Support for ORCID icon
\usepackage{orcidlink}



\begin{document}

% ---------------------------------------------------------------
% TODO REVIEW: Replace with your title
\title{\ours{}: Can Foundation Image Models be Agents for 3D World Models?\vspace{-10px}} 

% TODO REVIEW: If the paper title is too long for the running head, you can set
% an abbreviated paper title here. If not, comment out.
\titlerunning{\ours{}}

% TODO FINAL: Replace with your author list. 
% Include the authors' OCRID for the camera-ready version, if at all possible.
\author{Ziya Erko\c{c} \and
Angela Dai \and
Matthias Nie{\ss}ner \vspace{-5px}}

% TODO FINAL: Replace with an abbreviated list of authors.
\authorrunning{Z.~Erko\c{c} et al.}
% First names are abbreviated in the running head.
% If there are more than two authors, 'et al.' is used.

% TODO FINAL: Replace with your institution list.
\institute{Technical University of Munich \\ \vspace{5px}
\url{https://ziyaerkoc.com/worldagents} \vspace{-5px}}

\maketitle
\begin{figure}
    \centering
    \scalebox{1.0}{%
        % The minipage groups the image and caption into one block
        \begin{minipage}{0.88\textwidth}
    \includegraphics[width=\linewidth]{figures/3DAgent_Teaser.pdf}
    \caption{\textbf{\ours} employs 2D foundation models as agents in an iterative process to extract realistic, coherent 3D scenes from their learned distributions. 
    The scene generation process is orchestrated by a VLM agent (director) that generates prompts for new views to synthesize, fulfilled by an image generation model that creates these views, and verified by a VLM agent that assesses both 2D and 3D consistency of the new synthesized image(s).
    The output is a sci-fi lab. scene reconstructed with 3DGS, and can be visualized from novel views for exploration.}
    \vspace{-15px}
    \label{fig:teaser}
    \end{minipage}
    }
\end{figure}
\begin{abstract}
We introduce \textit{SegviGen}, a framework that repurposes native 3D generative models for 3D part segmentation. 
%
Existing pipelines either lift strong 2D priors into 3D via distillation or multi-view mask aggregation, 
%
often suffering from cross-view inconsistency and blurred boundaries, 
%
or explore native 3D discriminative segmentation, 
%
which typically requires large-scale annotated 3D data and substantial training resources.
%
In contrast, \textit{SegviGen} leverages the structured priors encoded in pretrained 3D generative model to induce segmentation through distinctive part colorization, establishing a novel and efficient framework for part segmentation.
%
Specifically, \textit{SegviGen} encodes a 3D asset and predicts part-indicative colors on active voxels of a geometry-aligned reconstruction. 
%
It supports interactive part segmentation, full segmentation, and full segmentation with 2D guidance in a unified framework.
%
% Extensive experiments show that \textit{SegviGen} outperforms prior state of the art by approximately 40\% on interactive part segmentation and 15\% on full segmentation, while using only 0.32\% of the training data.
Extensive experiments show that \textbf{\textit{SegviGen} improves over the prior state of the art by 40\% on interactive part segmentation and by 15\% on full segmentation, while using only 0.32\% of the labeled training data.}
%
It demonstrates that pretrained 3D generative priors transfer effectively to 3D part segmentation, enabling strong performance with limited supervision.
\end{abstract}
% teaser % lin
% 3D 打印视频
% 渲染视频
\section{Introduction}

Part segmentation provides explicit part-level structures of 3D assets, serving as a core primitive for 3D content creation pipelines and offering fundamental 3D perception capabilities for spatial intelligence.
%
It enables a wide range of downstream applications, including part-level editing, animation rigging, and industrial uses such as 3D printing.
%
However, existing methods often fall short in segmentation quality, producing erroneous regions and imprecise boundaries that limit their practical usability.


To this end, one line of work attempts to transfer the comprehensive 2D segmentation priors to 3D via 2D-to-3D lifting.
%
Methods such as SAMPart3D~\cite{sampart3d} optimize 3D segmentation via 2D-to-3D distillation, 
but incur substantial computational and time overhead, and often yield blurry boundaries.
%
In parallel, another set of methods~\cite{segment3d,pointsam,geosam2} applies SAM~\cite{sam,sam2,sam3} to obtain 2D masks of multi-view projected images, which are then back-projected and fused into 3D masks.
%
However, these multi-view pipelines incur substantial runtime overhead, are sensitive to view coverage, and the back-projection and fusion step often introduces cross-view inconsistencies and imprecise boundaries.

% su du kuai
Recently, another line of work~\cite{p3sam,partsam} moves toward native 3D part segmentation so as to remedy the inherent shortcomings of the aforementioned methods that leverage 2D segmentation priors.
%
These methods predict segmentation parts directly in the native 3D space, explicitly enforce semantic and structural consistency, and are more efficient at inference.
%
However, it is a typical requirement to collect large-scale training datasets with curated 3D part annotations, 
where fine-grained annotations are costly and inconsistent across sources in granularity, hierarchy, and boundary definitions.


%
% 总结，2D的问题，mismatch，3D的问题，然后本质promising是纯三维，41行，我们发现能够物体结构，外观这里，分割，prior
% Overall, existing approaches face a fundamental trade-off between annotation and training cost and the structural artifacts caused by lifting 2D cues into 3D.
% %
% We attribute these artifacts to the inherent mismatch between 2D priors and native 3D structure,
% %
% motivating us to instead exploit a native 3D structural prior in place of lifted 2D cues.
In summary, the first line of methods suffers from a mismatch between 2D priors and 3D structure, while the second relies on costly training from scratch.
%
Therefore, a more promising approach is to leverage a prior model that encodes both 3D structure and texture to perform segmentation.
%
In particular, 3D generative models trained on large-scale unannotated 3D textured assets internalize rich part-level structure and texture patterns, providing a strong 3D prior over geometry and appearance. 
%
Such priors encourage part segmentation with sharper boundaries, while reducing reliance on dense part annotations and extensive task-specific training.
%
This motivates us to ask:
\textit{How can 3D generative priors be effectively transferred to part-level 3D segmentation to improve quality and data efficiency?}


Motivated by this perspective, we propose \textit{SegviGen}, a generative framework for 3D part segmentation that leverages the rich 3D structural and textural knowledge encoded in large-scale 3D generative models.
%
Specifically, we formulate part segmentation as a colorization task that fully exploits the capacity of 3D generative models.
%
\textit{SegviGen} encodes the input 3D asset into a latent representation and uses it, together with the task embedding and query points, to condition the denoising process.
%
The model is trained to predict part-indicative colors, along with reconstructing the underlying geometry.
%
This formulation naturally accommodates additional conditioning signals, enabling \textit{SegviGen} to flexibly support interactive part segmentation, full segmentation, and 2D segmentation map–guided full segmentation under a unified architecture. 
%
Notably, while the first two tasks are common settings, 2D segmentation map–guided full segmentation is uniquely enabled by \textit{SegviGen}, supporting arbitrary part granularity and more precise segmentation that is critical for industrial applications.
% 写清楚哪几个setting


Qualitative and quantitative results show that \textit{SegviGen} consistently surpasses the prior state of the art, P3-SAM~\cite{p3sam}, while using only 0.32\% of the labeled training data. % 写清楚啥方法。
%
On interactive part segmentation, it achieves the best performance across all metrics on PartObjaverse-Tiny~\cite{partobjaverse_tiny} and PartNeXT~\cite{partnext}, with a 40\% gain in IoU@1, an important metric that reflects the model’s single-click accuracy. 
% 体现这个指标的挑战性 
%
On full segmentation without guidance, \textit{SegviGen} outperforms the best baseline by 15\% in overall IoU, averaged across datasets.
%
Our main contributions are summarized as follows:
\begin{itemize}[leftmargin=*, topsep=2pt, itemsep=1pt, parsep=0pt, partopsep=0pt]
    \item We propose \textit{SegviGen}, a unified multi-task framework for 3D part segmentation that effectively exploits the structural and textural priors encoded in pretrained 3D generative models, enabling accurate and efficient segmentation.
    \item We reformulate 3D segmentation as part-wise colorization, where \textit{SegviGen} predicts the colors of actiave voxel as part labels in a single generative process.
    \item Extensive experiments show that \textit{SegviGen} outperforms the prior state of the art by 40\% on interactive part segmentation and 15\% on full segmentation, using only 0.32\% of the labeled training data, highlighting the effectiveness of transferring 3D generative priors to part segmentation.
\end{itemize}


\input{sections/2_related}
\section{Methodology}

We propose \textit{SegviGen}, a unified multi-task framework for 3D part segmentation that supports three practical settings: interactive part-segmentation, full segmentation, and full segmentation with 2D guidance.
%
To leverage the prior knowledge encoded in a pretrained 3D generative model, we cast 3D segmentation as a colorization problem.
%
Specifically, we encode the input 3D asset into a compact latent that conditions generation, and optionally augment it with user interactions or a 2D segmentation map.
%
Conditioned on these inputs, the model reconstructs the 3D asset while predicting colors for active voxels in the structured 3D representation, where each color corresponds to an individual part, yielding the final segmentation.
%
Below, we begin by describing the underlying 3D generative model (Sec.~\ref{sec:3d_generative_model}), followed by our task reformulation (Sec.~\ref{sec:task_define}), and then detail the overall pipeline (Sec.~\ref{sec:pipeline}).


\subsection{Preliminary: Structured 3D Generative Model}
\label{sec:3d_generative_model}

Recent work~\cite{trellis2} organizes each textured 3D asset into a sparse set of active voxels on a regular grid, where every active voxel stores geometry and texture features aligned in 3D.
%.
This formulation leverages a flexible dual-grid construction to robustly handle arbitrary topology while encoding physically-based material attributes jointly with geometry for faithful appearance modeling.
%
Given the sparse omni-voxel representation, a Sparse Compression VAE (SC-VAE) maps each voxelized asset feature tensor $\mathbf{x}$ to a compact structured latent $\mathbf{z}_1=E_{\phi}(\mathbf{x})$ and reconstructs it via $\hat{\mathbf{x}}=D_{\theta}(\mathbf{z}_1)$, yielding an expressive yet highly compressed 3D latent space.
%
On top of these latents, a conditional flow-matching generator learns a time-dependent vector field $\mathbf{v}_{\psi}(\mathbf{z}_t,t,\mathbf{c})$ under conditioning $\mathbf{c}$ by matching the constant velocity along linear interpolants:
\begin{equation}
\mathbf{z}_0\sim\mathcal{N}(\mathbf{0},\mathbf{I}),\quad t\sim\mathcal{U}(0,1),\quad \mathbf{z}_t=(1-t)\mathbf{z}_0+t\mathbf{z}_1,
\end{equation}
\begin{equation}
\mathcal{L}_{\mathrm{cfm}}
=\mathbb{E}\Big[\big\|\mathbf{v}_{\psi}(\mathbf{z}_t,t,\mathbf{c})-(\mathbf{z}_1-\mathbf{z}_0)\big\|_2^2\Big].
\end{equation}
%
This latent generative pipeline enables efficient synthesis of geometry- and texture-consistent 3D assets, 
and the resulting structured latents capture rich joint statistics of shape and appearance, providing a strong transferable prior for fine-grained 3D part segmentation.


\subsection{Task Reformulation and I/O Representation}
\label{sec:task_define}

We reformulate 3D part segmentation as a color prediction problem in a structured 3D representation.
%
This choice matches our base 3D generative model~\cite{trellis2}, which jointly parameterizes geometry together with appearance attributes such as color, material properties, and roughness in a unified representation.
%
To maximize reuse of the pretrained generative prior, 
we avoid introducing an additional segmentation-specific attribute channel, which would increase modeling and optimization complexity, and instead express segmentation targets directly in color space, the most visually intuitive attribute.
We consider three task settings with consistent input and output formats.

%
\textbf{Interactive part-segmentation} is formulated as binary part extraction: given 3D points indicating a target part, we supervise the model to color the selected part in white and the remaining regions in black.
%
\textbf{Full segmentation} targets multi-part decomposition: we assign each part a distinct color from a randomly sampled color palette and supervise voxel colors accordingly.
%
Importantly, correctness is defined up to a permutation of colors within each object; any one-to-one assignment between predicted colors and parts is considered valid.
%
To reduce sensitivity to particular color choices, we use $K{=}10$ independently sampled palettes per shape, providing multiple colorizations for the same underlying partition.
%
\textbf{Full segmentation with 2D guidance} additionally conditions the model on a rendered 2D segmentation map: we first colorize the 3D parts and render the corresponding 2D segmentation map, and we then train the model to generate 3D voxel colors that are consistent with the color assignments in the 2D guidance.
%
Overall, this formulation preserves a unified model interface across settings, enabling a consistent architecture and training pipeline.


\subsection{Unified Multi-Task 3D Part Segmentation}
\label{sec:pipeline}

\paragraph{Overall framework.}
To fully leverage pretrained 3D generative models, we cast 3D part segmentation as a conditional part-wise colorization task in 3D latent space.
%
Given an input asset $X$, a pretrained 3D VAE encoder $E(\cdot)$ produces a encoded latent $z = E(X)$, 
which helps specify the active voxel support and anchors generation to the underlying shape.
%
For each task, we construct a part-wise colorized target and encode it into the same latent space to obtain $y$, following the task-specific scheme in Sec.~\ref{sec:task_define}.
%
We then sample $\epsilon \sim \mathcal{N}(0,I)$ and $t\sim\mathcal{U}(0,1)$ to form a noisy interpolation
\begin{equation}
y_t \;=\; (1-t)\,y \;+\; t\,\epsilon .
\end{equation}
%
A pretrained DiT-based backbone is fine-tuned to predict the noise residual conditioned on the noisy input $y_t$, the geometry latent $z$, the task condition $C$, and a learned task embedding $e_{\tau}$:
\begin{equation}
\hat{v}_\theta \;=\; f_\theta\!\left(y_t,\, z,\, C,\, e_{\tau},\, t\right).
\end{equation}
%
Training follows the conditional flow-matching objective
\begin{equation}
\mathcal{L}(\theta)
\;=\;
\mathbb{E}_{X,\tau,t,\epsilon}\Big[
w(t)\,\big\| \hat{v}_\theta \;-\; (\epsilon - y) \big\|_2^2
\Big],
\end{equation}
where $w(t)$ is an optional timestep weighting.



\begin{figure*}
  \centering
  \includegraphics[width=\linewidth]{img/comparison_point.pdf}
  \caption{
  Interactive part-segmentation results. 
  %
  We compare \textit{SegviGen} with existing representative baselines, including Point-SAM~\cite{pointsam} and P3-SAM~\cite{p3sam}. 
  %
  In the figure, yellow points denote user clicks, while the predicted target part is highlighted in red.
  %
  Leveraging priors from pretrained 3D generative models, \textit{SegviGen} achieves more accurate results with sharper boundaries than prior methods, while requiring substantially less training data.
  }
  \label{fig:comparison_point}
\end{figure*}


\paragraph{Condition Injection.}
We adopt task-specific conditioning designs while maintaining a unified interface across settings.
%
For interactive segmentation, user clicks in the UI provide an efficient and intuitive form of guidance.
%
In our framework, each click is encoded as a sparse point token comprising its 3D coordinates and an associated feature vector.
%
Since the 3D coordinates are already effectively encoded by RoPE within the attention layers, 
%
we omit the additional learnable input-level positional embedding used in prior designs~\cite{p3sam}.
%
Instead, all points share the same learnable feature vector $Q$ , which serves as the point token during both training and inference.
%
Given point coordinates $\mathbf{u}=\{u_i\}_{i=1}^{m}$ with $u_i\in\mathbb{R}^3$, we form point-condition tokens
\begin{equation}
Q \;=\; \big[\; \mathbf{q}(u_1),\ldots,\mathbf{q}(u_m) \;\big], 
\quad
\mathbf{q}(u_i) \;=\; \big[\, u_i \,;\, \mathbf{e}_p \,\big],
\end{equation}
where $\mathbf{e}_p$ is a shared learnable feature appended to every point token.
%
Conditioned on $Q$, the denoising model is instantiated as
\begin{equation}
\hat{v}_{\theta} \;=\; f_{\theta}\!\left(y_t,\, z,\, Q,\, e_{\tau},\, t\right).
\end{equation}
%
When the number of points is fewer than $10$, we pad the point tokens to a length of $10$
using zero coordinates and zero features.
%
To preserve a single unified model, we keep this interface for full segmentation and 2D-guided full segmentation by providing $10$ padded tokens with all-zero coordinates and features.


For full segmentation with 2D guidance, we additionally provide a user-specified 2D segmentation colorization as guidance, allowing direct control over the desired part granularity and label palette.
%
The guidance image is encoded into a sequence of conditioning tokens injected via cross-attention:
\begin{equation}
p \;=\; g_{\phi}(I_{\text{guide}}),
\end{equation}
where $g_{\phi}(\cdot)$ denotes an image encoder.
%
In this setting, denoising is conditioned on both the padded point-token interface $Q_{0}$ and the image guidance tokens $p$:
\begin{equation}
\hat{v}_{\theta} \;=\; f_{\theta}\!\left(y_t,\, z,\, (Q_{0},\, p),\, e_{\tau},\, t\right).
\end{equation}

\paragraph{Task Embedding.}
To improve multi-task generalization within a single model, task identity is encoded as a continuous embedding and injected alongside the timestep signal.
%
Let $\tau\in\{1,\dots,T\}$ denote the task index.
%
A sinusoidal encoding is first computed from $\tau$,
\begin{equation}
s_{\tau} \;=\; \mathrm{PE}(\tau)\in\mathbb{R}^{d_f},
\end{equation}
where $\mathrm{PE}(\cdot)$ follows the standard sinusoidal scheme.
%
A lightweight MLP then maps $s_{\tau}$ to the task embedding
\begin{equation}
e_{\tau} \;=\; \mathrm{MLP}_{\psi}(s_{\tau}) \in \mathbb{R}^{d}.
\end{equation}
In parallel, the timestep $t$ is embedded as $e_t\in\mathbb{R}^{d}$.
The final modulation vector used by DiT backbone is obtained by additive fusion,
\begin{equation}
m \;=\; e_t \;+\; e_{\tau},
\qquad
\hat{v}_{\theta} \;=\; f_{\theta}\!\left(y_t,\, z,\, C,\, m\right),
\end{equation}
where $m$ conditions the adaptive layers to jointly encode diffusion progress and task semantics.
%
During training, samples from different tasks are interleaved and supervised with their corresponding $\tau$, 
encouraging the shared backbone to learn task-discriminative behaviors while preserving a unified parameterization.











\input{sections/4_experiments}
\section{Conclusion}

This paper introduces \textit{SegviGen}, a framework that repurposes pretrained 3D generative models for 3D part segmentation. 
%
In contrast to 2D-to-3D lifting methods that often suffer from cross-view inconsistency and blurred boundaries, 
and native 3D discriminative approaches that require large-scale part annotations and heavy training, 
\textit{SegviGen} transfers generative priors to deliver accurate and globally coherent segmentations with limited supervision. 
%
It reformulates segmentation as part-wise colorization, jointly reconstructing geometry and predicting part-indicative colors, 
and supports multiple task settings via flexible conditioning. 
%
Experiments on interactive and full segmentation benchmarks show consistent improvements over prior methods, underscoring the effectiveness and data efficiency of 3D generative priors for 3D part segmentation.

\begin{figure*}[t]
  \centering

  % Row 1
  \begin{minipage}[t]{0.49\linewidth}
    \centering
    \includegraphics[width=\linewidth]{img/w1.png}\\[-0.2em]
    % {\small (a)}
  \end{minipage}
  \hspace{0.01\linewidth}
  \begin{minipage}[t]{0.49\linewidth}
    \centering
    \includegraphics[width=\linewidth]{img/w2.png}\\[-0.2em]
    % {\small (b)}
  \end{minipage}

  \vspace{0.01\linewidth}

  % Row 2
  \begin{minipage}[t]{0.49\linewidth}
    \centering
    \includegraphics[width=\linewidth]{img/w3.png}\\[-0.2em]
    % {\small (c)}
  \end{minipage}
  \hspace{0.01\linewidth}
  \begin{minipage}[t]{0.49\linewidth}
    \centering
    \includegraphics[width=\linewidth]{img/w4.png}\\[-0.2em]
    % {\small (d)}
  \end{minipage}

  \caption{Interactive demo. Users specify clicks on the 3D asset to perform interactive part segmentation, and can adjust visualization settings.}
  \label{fig:web_screenshot}
\end{figure*}

\begin{figure*}
  \centering
  \includegraphics[width=\linewidth]{img/comparison_more_result_v2.pdf}
  \caption{More qualitative comparisons for full segmentation.}
  \label{fig:comparison_more_result}
\end{figure*}



% ---- Bibliography ----
%
% BibTeX users should specify bibliography style 'splncs04'.
% References will then be sorted and formatted in the correct style.
%
\bibliographystyle{splncs04}
\bibliography{main}
\clearpage
\input{sections/6_supp}

\end{document}
